\documentclass[letterpaper,12pt]{article}

% Librerías a utilizar
\usepackage[utf8]{inputenc}	% Codificación 
\usepackage[T1]{fontenc}
\usepackage[style=apa, backend=biber, language=spanish]{biblatex}
\usepackage{graphicx}		% Igenes
\usepackage{indentfirst}		% Sangría
\usepackage{amsmath, amsfonts, amssymb}	% Figuras matemáticas
\usepackage{xurl}    % URL
\usepackage{hyperref}
\usepackage{csquotes}
\usepackage{booktabs}
\usepackage{placeins}
\usepackage{tabularx}
\usepackage[spanish,es-tabla]{babel}	% Idioma
\addbibresource{bibliografia.bib}

\usepackage[left=3cm,right=3cm,top=2.54cm,bottom=2.54cm]{geometry}
\usepackage{multirow}

\setlength{\parindent}{2cm}	% Sangría en los párrafos
\renewcommand{\baselinestretch}{1.5}	% Interlineado

% Renombrar ciertos títulos del texto
\renewcommand\spanishcontentsname{Tabla de contenidos}
\renewcommand\spanishrefname{Bibliografía}

% Inicio del documento
\begin{document}

\newpage
\vspace*{-.5cm}
% Logo institucional
\begin{picture}(18,4)(0,30)
	\put(370,-20){
\includegraphics[scale=0.16]{./images/Logo_Usach.pdf}}
\end{picture}

\sloppy
\thispagestyle{empty}
\vspace*{-1.6cm}

% Datos institucionales
\begin{center}
\vspace{0.5cm}
	{\bf \mbox{\Large UNIVERSIDAD DE SANTIAGO DE CHILE}}\\
	{\bf \mbox{\large FACULTAD DE INGENIERÍA}}\\
	{\bf \mbox{\large Departamento de Ingeniería Informática}}\\
\end{center}

	\vspace{5cm}
 
	%Título del trabajo
	\begin{center}
	\Large
		\textbf{Coord-Agent: Agente que apoya la coordinación de actividades grupales}\\
	\end{center}
	
	% Datos personales
	\vspace*{5cm}
	\begin{flushright}
		\begin{tabular}[t]{l l}
			Integrantes:&Daniel Eguiluz \\
                        &Sergio Espinoza\\
						&Matías Hernández\\
                        &Nicolás Sánchez\\
			Asignatura: &TAVI - Diurno\\
			Profesor:   &Daniel Gacitúa

		\end{tabular}
	\end{flushright}
	\begin{center}
		\vspace{1.5cm}
		% Fecha
        Santiago - Chile\\
		\today
	\end{center}

\newpage

\include{chapters/0-resumen}

\normalsize
\tableofcontents
\thispagestyle{empty}

\newpage
\renewcommand{\thepage}{\arabic{page}}
\setcounter{page}{1}

% Capítulos
\section{Introducción} \label{sec:introduccion}

a


\section{Definición del proyecto}

\subsection{Identificación del problema u oportunidad}

b

\subsection{Justificación del propósito}

c

\subsection{Planteamiento de los objetivos del proyecto}

d
\section{Estudio de factibilidad}

\subsection{Análisis técnico}

e

\subsection{Análisis económico}

f

\subsection{Evaluación de riesgos y limitaciones}

g
\section{Metodología de desarrollo}

\subsection{Selección de una metodología ágil}

Se a decidido que para el desarrollo de este proyecto se utilizará la metodología Scrum, esta metodología altamente conocida por su enfoque iterativo incremental promueve la adaptabilidad y la entrega continua, de esto se desligan distintas razónes de porque fue seleccionada:

\begin{enumerate}
    
    \item \textbf{Iteraciones cortas (Sprints):} Los ciclos cortos de Scrum son la clave para este proyecto, el poder contar con versiones funcionales desde etapas tempranas permitirá medir el avance y el valor que esta entregando la solución casi desde el principio, y así evaluar la viabilidad de posibles incrementos.

    \item \textbf{Retroalimentación temprana:} Siguiendo con lo del punto anterior, estas versiones funcionales casi al inicio del ciclo de vida del proyecto permitirá realizar pruebas que permitan detectar de forma oportuna errores o desviaciones, mejorando la calidad del producto final.
    
    \item \textbf{Alta adaptabilidad a cambios:} La posibilidad mencionada anteriormente de recibir retroalimentación continua sobre el producto permitirá adaptarse a cualquier cambio que sea necesario que fuera imprevisto durante la planificación inicial. Esto es clave teniendo en cuenta que es el primer proyecto basado en LLM de los integrantes del grupo, por lo que se podrían pasar cosas por alto en una planificación inicial por mera falta de experiencia en el área.
\end{enumerate}


\subsection{Planificación de iteraciones y definición de roles en el equipo}

\subsubsection{Plan de iteraciones}

a

\subsubsection{Roles del equipo y tareas asociadas}

\begin{enumerate}
    \item Nicolás Sánchez - Scrum Master: Encargado de facilitar la adopción de la metodología Scrum por parte del equipo y asegurando el cumplimiento de las prácticas y entregas asociadas a dicha metodología, además de coordinar el flujo de trabajo y reuniones grupales. 
    
    \item Matías Hernández - Product Owner: Encargado de definir los requerimientos del producto, validar el cumplimiento de este hacía las necesidades de los posibles usuarios desde la toma de decisiones sobre los pasos a seguir para obtener un buen producto, hasta que este sea desplegado para los usuarios finales.
    
    \item Desarrollador backend - Sergio Espinoza: Encargado del diseño e implementación del backend, base de datos e integración de módulos.
    
    \item Desarrollador frontend - Daniel Eguiluz: Encargado del diseño e implementación de la interfaz para los usuarios y para el consumo del backend.
\end{enumerate}

Para la selección de estos roles y respectivas tareas se tomaron en cuenta los trabajos previos realizados por cada miembro, de forma que vaya acorde a sus áreas de especialización y así asegurar el cumplimiento de los objetivos del proyecto.


\subsection{Diagrama de usos éticos de la IA (MyAI Compass)}

a

\subsection{Repositorio de control de versiones}

Para el control de versiones del código y los informes del proyecto se hará uso de github, específicamente del repositorio presente en el siguiente enlace: \url{https://github.com/Ch3chS/TAVI-Charlie}.
\section{Diseño arquitectural}

\subsection{Definición de la arquitectura de software}

La solución que se ha implementado cuenta con un patrón cliente web + BFF + servicios internos. El frontend, desarrollado en React consume rutas del backend FastAPI. El backend actua como BFF se encarga de la extracción del LLM, la fusión de los datos dentro de la sesión, el cálculo de opciones, la persistencia y las respuestas para la interfaz de usuario (canal). Para la creación de la aplicación el archivo \texttt{main.py} se encarga del CORS; \texttt{routes.py} concentra las interfaces; \texttt{session\_service.py} se encarga de las sesiones; \texttt{decision\_engine.py} se encarga de la toma de decisión sobre la mejor sesión; \texttt{json\_repository.py} se encarga de la persistencia de los datos.

Además, se considera una capa de integración para canales de comunicación. En el estado actual del proyecto, el sistema cuenta con un canal simulado para las conversasiones grupales. Para el desarrollo del MVP se contempla la implementación de al menos una integración real, priorizando Discord debido a su simplicidad de integración. La capa quedara diseñada para permitir futuras integraciones ya sea con WhatsApp, Slack u otras plataformas.
Finalmente, se incorporará un servicio de calendario que permitira consolidar e integrar la decisión confirmada en un evento calendarizable, inicialmente en un formato .ics.

El flujo general quedaria asi: mensaje del usuario $\rightarrow$ canal web o externo $\rightarrow$ BFF $\rightarrow$ servicio LLM extractor $\rightarrow$ actualización de sesión $\rightarrow$ motor de decisión $\rightarrow$ respuesta al usuario $\rightarrow$ confirmación $\rightarrow$ generación de evento de calendario.

\begin{figure}[htbp]
    \centering
    \includegraphics[width=0.9\textwidth]{Arquitectura_General.png}
    \caption{Diagrama Arquitectura general}
    \label{fig:arquitectura_general}
\end{figure}

\subsection{Identificación de la integración del LLM en la solución}
El LLM se integra como un componente dentro del sercio \textit{LLM Service}. Es importante que la función no es tomar una decisión final, sino transformar lo mensajes en lenguaje natural a datos estructurados en un formato especifico para que puedan ser utilizados el resto del sistema.

Por ejemplo, si llegase un mensaje como ``Nicolás puede el lunes en la tarde y Matías puede desde las 16'', el LLM debe extraer participantes, días y horarios disponibles. Esta información es devuelta en un formato estructurado en JSON, para que el backend lo procese y lo fusione con el resto de la información de la sesión.

La decisión final depende del servicio \textit{Decision Engine}, el servicio se encarga de calcular de manera deterministica la mejor opcion evaluando la cobertura de los participantes y generando una explicación para cada alternativa generada.

La integración considera mecanismos basicos de seguridad como prompts restrictivos, validación de la salida, correccion manual de los datos reconocidos por el sistema, uso de fallback en caso de falllas del modelo y la confirmación explicita del usuario a la hora de la toma  de decisión final.

\begin{figure}[htbp]
    \centering
    \includegraphics[width=1.1\textwidth]{LLM_en_solucion.png}
    \caption{Flujo de integración del LLM en el proceso de coordinación.}
    \label{fig:flujo_LLM}
\end{figure}

\subsection{Diagramas UML y especificaciones técnicas}

Para el diseño del sistema se consideran tres vistas principales: una vista de componentes, una vista de flujo y una vista de clases. La vista de componentes permite observar la separación entre frontend, backend, servicios internos e integraciones externas. La vista de flujo permite describir cómo un mensaje en lenguaje natural se transforma en una decisión grupal. Finalmente, la vista de clases permite representar las entidades principales utilizadas por el sistema.

Las clases principales del modelo son \texttt{Session}, \texttt{Participant}, \texttt{TimeSlot}, \texttt{TimeOption}, \texttt{AvailabilityCell}, \texttt{ChannelMessage} y \texttt{ChannelConfig}. La clase central es \texttt{Session}, ya que representa una coordinación específica. Esta contiene participantes, mensajes, disponibilidades, opciones calculadas, matriz de disponibilidad y, en caso de confirmación, una opción seleccionada y un resumen final de la decisión.

En este diseño, \texttt{Participant} solo corresponde a una persona detectada en una sesión, no a una cuenta global. Si se requiere en un futuro deberia crearse la clase \texttt{User} para soportar esto.

A nivel técnico, el MVP utiliza React y TypeScript para el frontend, FastAPI y Python para el backend, comunicación mediante HTTP/REST, persistencia local en JSON, extracción mediante LLM configurable y un motor determinístico para el cálculo de opciones. Además, se contempla una integración mínima con un canal de comunicación externo y una integración con calendario mediante generación de archivo \texttt{.ics}.

\begin{figure}[htbp]
    \centering
    \includegraphics[width=1.1\textwidth]{Clases.png}
    \caption{Diagrama de clases principales del sistema.}
    \label{fig:clases}
\end{figure}

\section{Mini-demo funcional}

\subsubsection{Despliegue de una instancia mínima de un LLM}

m
\section{Conclusiones}

En esta primera entrega se ha establecido la estructura base del proyecto. A partir de la identificación de una problemática real y frecuente en entornos organizacionales (la ineficiencia en la coordinación grupal) se definió una solución centrada en un agente con un LLM como núcleo. El estudio de factibilidad evidencia que el proyecto es técnica y económicamente realizable dentro del plazo de un semestre, con un stack tecnológico relativamente accesible. La adopción de Scrum como metodología de desarrollo, junto con la definición clara de roles y un diseño arquitectural modular, sienta las condiciones para un avance iterativo y controlado hacia el MVP. Las siguientes etapas del proyecto se enfocarán en la implementación progresiva de los sprints definidos, con miras a demostrar en la Evaluación 2 un avance funcional de al menos el 50\% de los objetivos comprometidos, tal como se pide.


\clearpage
\addcontentsline{toc}{section}{Bibliografía}
\printbibliography

\include{chapters/5-Anexo}

\end{document}